% titel: Blatt 4, A1: Konvergenz von Folgen und Reihen

% Teil I: Werkzeuge, kein Lösungsweg.

Konvergenz einer Wurzelfolge mit explizitem $N$ (je nach $\varepsilon$),
Nichtkonvergenz einer alternierenden Folge und eine konkrete Reihe.

\begin{werkzeugkasten}[Definition Konvergenz]
$a_n\to a$: zu jedem $\varepsilon>0$ existiert $N$ mit $|a_n-a|<\varepsilon$
für alle $n>N$. Für $a_n=\frac1{\sqrt{n^3+2}}\to0$ löst man
\begin{equation}
  \frac1{\sqrt{n^3+2}}<\varepsilon\ \Longleftrightarrow\ n^3+2>\frac1{\varepsilon^2}
\end{equation}
und wählt $N$ entsprechend der linken Seite von $\varepsilon$ (und damit
auch für $\varepsilon=10^{-3}$ konkret $N=99$).
\end{werkzeugkasten}

\begin{werkzeugkasten}[Keine Konvergenz bei alternierenden Folgen]
Eine Folge $b_n=(-1)^n$ konvergiert nicht: die Teilfolge
$b_{2k}=1$ und $b_{2k+1}=-1$ hätten verschiedene Grenzwerte, was eine
eindeutiger Konvergenz ausschließt. Eine konvergente Folge hat genau einen
Grenzwert.
\end{werkzeugkasten}

\begin{werkzeugkasten}[Geometrische Reihe]
\begin{equation}
  \sum_{k=0}^{\infty}q^k=\frac1{1-q}\quad(|q|<1).
\end{equation}
Für $\sum_{k=2}^{\infty}(-1)^k\frac{3^k}{4^{k+1}}$ klammert man
$\frac14$ aus und summiert die geometrische Reihe mit
$q=-\frac34$; dabei die fehlenden Glieder $k=0,1$ wieder abziehen.
\end{werkzeugkasten}

\begin{achtung}[Darauf achten]
\begin{itemize}
  \item Das $N$ muss zu $\varepsilon$ passen: aus
        $n^3+2>\frac1{\varepsilon^2}$ folgt ein $N(\varepsilon)$; für das
        konkrete $\varepsilon=10^{-3}$ gibt das $N\ge100$.
  \item ``$n>N$'' und ``$n\ge N$'' nicht verwechseln -- entsprechend wählen.
  \item $(-1)^n$ hat keine Grenzwert-Eigenschaft; der Vergleich zweier
        Teilfolgen zeigt die Nichtkonvergenz.
  \item Bei der Reihe zuerst die Indizes $k=0,1$ vom Rest trennen, dann
        geometrisch aufsummieren.
\end{itemize}
\end{achtung}
